\section{Related work}

This paper sits at the intersection of software contracts, static-versus-dynamic checking, and typed data-pipeline boundaries. The closest primary sources for the implementation itself are the Scala~3 reflection documentation and the Spark schema-comparator APIs, because the artifact is built directly on those mechanisms~\cite{scala3_reflection,spark_datatype}. These are mechanism sources, not related-work contributions in the research sense.

The closest implementation-level neighbors are typed-Dataset libraries. \texttt{frameless} and \texttt{iskra} provide typed Spark layers that carry schema information through Dataset operations, but they do so by asking the user to replace raw \texttt{DataFrame} code with typed wrappers across the job~\cite{frameless,iskra}. This artifact targets a narrower seam: producer-to-contract conformance is required only at the checked boundary, and the rest of the job may remain in raw Spark code.

Record-to-record derivation libraries solve a different adjacent problem. Chimney, and earlier Scala generic-derivation work built on \texttt{shapeless}, derive field-to-field transformations with rich compile-time diagnostics~\cite{chimney,shapeless}. That macro technique is close in spirit to the compile-time side of this artifact, but the goal is different: Chimney proves transformation existence, whereas this artifact proves structural conformance under an explicit policy family and couples that witness to a Spark sink-boundary pin.

Table-level schema enforcement is also close, but at a different layer. Delta Lake and Apache Iceberg validate writes against stored table metadata and support schema evolution at that boundary~\cite{delta_lake,apache_iceberg}. Those systems are valuable runtime seatbelts, but they operate after job submission and against stored schemas rather than against declared Scala producer and contract types. The artifact here is therefore complementary rather than competing.

The policy names \texttt{Backward} and \texttt{Forward} deliberately reuse vocabulary that is familiar from Avro and schema-registry compatibility systems~\cite{avro_spec,avro_compatibility}. In this paper, however, the relation is structural over Scala product types and Spark schemas. It is not a value-level compatibility proof for serialized messages or registered schemas.

The broader contract literature provides the right contrast. Empirical work such as \emph{Contracts in Practice} shows that contract disciplines remain relevant in real software, which supports the motivation for making contract intent explicit~\cite{contracts_in_practice}. More expressive systems such as \emph{Effectful Software Contracts} and \emph{Trace Contracts} show how much richer the contract space becomes when effects, histories, or traces are part of the property being checked~\cite{effectful_software_contracts,trace_contracts}. By comparison, this artifact is deliberately narrower: it checks structural compatibility at a typed boundary.

There is also a useful static-versus-dynamic comparison point. \emph{Dynamic Contract Analysis for Parallel Programming Models} treats dynamic contracts as complementary to static analysis rather than as a complete replacement~\cite{dynamic_contract_analysis}. That lesson aligns with this paper's design choice: compile-time proof and runtime checking do different jobs and should be combined at the boundary instead of being presented as competitors.

Finally, work outside the JVM ecosystem helps position the contribution. P3317R0 argues for contracts that can often be resolved at compile time so that runtime cost can be reduced or removed~\cite{compile_time_resolved_contracts}. The artifact here is not a direct analogue of that C++ proposal, but it shares the same directional idea: when part of a contract question is structural and visible to the compiler, deferring everything to runtime is unnecessary.

The novelty claim of this paper is therefore narrow. It is not a new general theory of contracts. It is a small, evidenced combination of compile-time structural proof and runtime boundary checking for a typed Scala and Spark setting.
